\documentclass[a4paper,10pt]{article}

\usepackage[utf8]{inputenc}
\usepackage[T1]{fontenc}
\usepackage[italian]{babel}
\usepackage[margin=2cm, top=1.8cm, bottom=1.8cm]{geometry}
\usepackage{titlesec}
\usepackage{enumitem}
\usepackage{hyperref}
\usepackage{xcolor}
\usepackage{fontawesome5}
\usepackage{array}
\usepackage{tabularx}
\usepackage{parskip}

% Colors
\definecolor{accent}{RGB}{80, 100, 200}
\definecolor{darkgray}{RGB}{50, 50, 50}
\definecolor{lightgray}{RGB}{130, 130, 130}

% Hyperref setup
\hypersetup{
    colorlinks=true,
    urlcolor=accent,
    linkcolor=accent
}

% Section formatting
\titleformat{\section}
  {\large\bfseries\color{darkgray}}
  {}{0em}{}
  [\vspace{0.3em}\color{lightgray}\hrule\vspace{0.4em}]

\titlespacing{\section}{0pt}{1.2em}{0.5em}

\pagestyle{empty}

\setlist[itemize]{noitemsep, topsep=2pt, leftmargin=1.2em}

% Custom commands
\newcommand{\entry}[4]{%
  \noindent\textbf{#1} \hfill \textcolor{accent}{#2}\\
  \textit{\textcolor{lightgray}{#3}}%
  \ifx&#4&\else\\\vspace{0.2em}#4\fi
  \vspace{0.5em}
}

\newcommand{\entrylist}[5]{%
  \noindent\textbf{#1} \hfill \textcolor{accent}{#2}\\
  \textit{\textcolor{lightgray}{#3}}%
  \ifx&#4&\else\\\vspace{0.2em}#4\fi
  \vspace{0.3em}
  \begin{itemize}#5\end{itemize}
  \vspace{0.3em}
}

\newcommand{\langentry}[3]{%
  \noindent\textbf{#1} — #2%
  \ifx&#3&\else\\\textcolor{lightgray}{\small #3}\fi
  \vspace{0.4em}
}

\newcommand{\skillgroup}[2]{%
  \noindent\textbf{#1:} #2\par\vspace{0.35em}
}

%--------------------------
\begin{document}

%--- HEADER ---%
\begin{center}
  {\Huge\bfseries\color{darkgray} Mattia Giovannini}\\[0.4em]
  {\large\color{accent} Studente Magistrale in Artificial Intelligence | Laureato in Ingegneria Informatica}\\[0.7em]
  \textcolor{lightgray}{
    \faPhone\ +39 XXXXXXXXXX
    \quad\textbullet\quad
    \faEnvelope\ \href{mailto:mattia.giovannini02@gmail.com}{mattia.giovannini02@gmail.com}
  }
\end{center}

\vspace{0.5em}

%--- PROFILE ---%
\section{Profilo}

Laureato in Ingegneria Informatica presso l’Università di Bologna e ammesso al corso di Laurea Magistrale in Artificial Intelligence presso lo stesso ateneo, con un forte interesse per l’intelligenza artificiale, il machine learning e l’analisi dei dati. Il mio percorso mi ha fornito solide basi in algoritmi, statistica e ingegneria del software, che ho applicato in un contesto di ricerca durante un semestre Erasmus in Francia. Sono motivato dalla sfida di progettare sistemi intelligenti per la risoluzione di problemi reali e pronto ad approfondire le mie competenze durante la laurea magistrale.

%--- EDUCATION ---%
\section{Formazione}

\entry{Laurea Magistrale in Artificial Intelligence}{2026 -- 2028 (Ammesso)}{Università di Bologna}
{Ammesso al corso di Laurea Magistrale a ciclo biennale per l'A.A. 2026/2027.}

\entrylist{Laurea Triennale in Ingegneria Informatica}{2021 -- Luglio 2026}{Università di Bologna}{Voto di laurea: 94/110}{
  \item Corsi rilevanti: Analisi Matematica, Algebra Lineare e Geometria,
        Matematica Applicata (Probabilità e Statistica), Fondamenti di Informatica,
        Sistemi Operativi, Reti di Calcolatori, Sistemi Informativi,
        Ingegneria del Software, Tecnologie Web, Architettura degli Elaboratori
  \item \textbf{Tesi:} \textit{Rilevamento del Disturbo dello Spettro Autistico tramite dati di Eye Tracking con tecniche di Machine Learning} — applicazione di tecniche di classificazione e analisi dei segnali a registrazioni cliniche di eye-tracking per il supporto allo screening dell’ASD
}

\entry{Diploma di Liceo Scientifico}{2016 -- 2021}{Liceo Augusto Righi, Bologna}{}

%--- RESEARCH EXPERIENCE ---%
\section{Progetto di Studio Erasmus}

\entrylist{Esperienza Erasmus}{2025}{Université de Tours, Francia}{}{
  \item Sviluppo di un progetto sul \textbf{machine learning per l’analisi del disturbo dello spettro autistico (ASD)}
  \item Applicazione di tecniche di classificazione ed estrazione di caratteristiche a dati di eye-tracking
  \item Esperienza maturata in un contesto accademico internazionale
  \item Rafforzamento delle competenze in analisi dei dati, Python e metodologia scientifica
}

%--- TECHNICAL SKILLS ---%
\section{Competenze Tecniche}

\skillgroup{Linguaggi di Programmazione}{C, Java, Python, JavaScript, HTML/CSS, SQL, Bash}
\skillgroup{Machine Learning e Analisi dei Dati}{scikit-learn, NumPy, pandas, Matplotlib, Jupyter Notebook}
\skillgroup{Sistemi e Networking}{Linux, TCP/IP, programmazione socket, database relazionali}
\skillgroup{Strumenti di Sviluppo}{Git, \LaTeX, MATLAB, utilizzo di Docker e Podman}
\skillgroup{Aree di Interesse}{Machine Learning,
                               Analisi dei Dati,
                               Containerizzazione, Automazione di Sistema}

%--- PROJECT EXPERIENCE ---%
\section{Esperienze di Progetto}

\entrylist{BOOM GEN AI Challenge 2026 — VEM Sistemi}{2026}{Gruppo di Progetto GenAI}{}{
  \item Partecipazione allo sviluppo di un \textbf{assistente basato su GenAI} per l'analisi automatica di bandi di gara, finalizzato al supporto del processo di valutazione \textbf{GO/NO-GO}.
  \item Progettazione e implementazione di un modulo di \textbf{scoring} in \textbf{Python}, pensato per valutare la coerenza tra requisiti del bando e capacità aziendali.
  \item Sperimentazione di tecniche di \textbf{prompt engineering} e integrazione con servizi \textbf{Azure OpenAI} per l'estrazione, classificazione e arricchimento dei requisiti.
  \item Contributo all'analisi della pipeline di \textbf{parsing documentale} per la normalizzazione di bandi in formato non strutturato.
}

%--- OTHER EXPERIENCE ---%
\section{Altre Esperienze}

\entrylist{Webmaster ed Editor Studentesco}{2020 -- 2021}{Liceo Augusto Righi, Bologna}{}{
  \item Progettazione e gestione delle piattaforme studentesche dell’istituto (\textit{Righi Network} e \textit{Il Sottobanco}), supportando l’interazione digitale e la presenza online del giornale scolastico.
  \item Sviluppo di competenze nelle tecnologie web, nella gestione dei contenuti e nella collaborazione in team.
}

\entry{Educatore Volontario presso Centro Estivo}{2016 -- 2020}{Estate Ragazzi, Bologna}
{Gestione di gruppi, lavoro di squadra, comunicazione e responsabilità in attività educative}

%--- LANGUAGES ---%
\section{Lingue}

\langentry{Italiano}{Madrelingua}{}

\langentry{Inglese}{C1}
  {Duolingo English Test: 130/160 (C1) · Certificazione CLA C1, Università di Bologna · Cambridge First Certificate B2 (178/190)}

\langentry{Francese}{B1/B2}
  {Certificazione CLA B1, Università di Bologna · Corsi B2/C1 frequentati presso l’Université de Tours}

\end{document}